\documentclass[11pt,a4paper]{article}
\usepackage[utf8]{inputenc}
\usepackage[T1]{fontenc}
\usepackage[spanish]{babel}
\usepackage[margin=2.5cm]{geometry}
\usepackage{booktabs}
\usepackage{array}
\usepackage{listings}
\usepackage{xcolor}
\usepackage{enumitem}
\usepackage{parskip}
\usepackage{fancyhdr}
\usepackage{amsmath}
\usepackage{amssymb}

\definecolor{codegray}{RGB}{245,245,245}
\definecolor{codeblue}{RGB}{0,60,180}
\definecolor{codegreen}{RGB}{0,120,0}
\definecolor{codepurple}{RGB}{120,0,120}

\lstdefinestyle{sql}{
  language=SQL,
  backgroundcolor=\color{codegray},
  basicstyle=\ttfamily\small,
  keywordstyle=\color{codeblue}\bfseries,
  commentstyle=\color{codegreen},
  stringstyle=\color{codepurple},
  breaklines=true,
  frame=single,
  framerule=0.4pt,
  xleftmargin=8pt,
  xrightmargin=8pt,
  numbers=none,
  tabsize=2,
  showstringspaces=false,
  morekeywords={FETCH,FIRST,ROWS,ONLY,OFFSET,OVER,PARTITION,DENSE_RANK,ROW_NUMBER,NTILE,LAG,LEAD,FIRST_VALUE,LAST_VALUE,INTERSECT,EXCEPT,AVG,SUM}
}
\lstset{style=sql}

\pagestyle{fancy}
\fancyhf{}
\rhead{\textcolor{gray}{\small TD7: Ingeniería de Datos — UTDT}}
\lhead{\textcolor{gray}{\small Examen Modelo 1 — Bloque 3 (Máquina)}}
\rfoot{\thepage}
\renewcommand{\headrulewidth}{0.4pt}

\setlist[enumerate,1]{label=\textbf{\arabic*.}}
\setlist[enumerate,2]{label=\alph*.}

\begin{document}

\begin{center}
  {\LARGE\bfseries TD7 — Ingeniería de Datos}\\[6pt]
  {\large Universidad Torcuato Di Tella}\\[4pt]
  {\Large\bfseries Examen — Modelo 1}\\[2pt]
  {\large\bfseries Bloque 3: Lenguajes de Consulta \textnormal{---} \textit{Parte en máquina}}\\[4pt]
  {\normalsize Duración: 1 hora \quad|\quad Puntaje: 3 puntos \quad|\quad Umbral de aprobación: 1{,}5\,/\,3}
\end{center}
\noindent\rule{\linewidth}{0.8pt}
\vspace{2pt}

\noindent\textbf{Instrucciones:} Escriba sus respuestas y envíelas por mail al finalizar la hora. Puede consultar la web para verificar \textbf{sintaxis SQL}. Está \textbf{prohibido} el uso de asistentes de IA (ChatGPT, Claude, Copilot, etc.), comunicarse con otras personas por cualquier medio, y compartir documentos en tiempo real (Google Docs, etc.).

\vspace{6pt}
\noindent\rule{\linewidth}{0.4pt}

% ─────────────────────────────────────────────────────────────
\subsection*{Esquema de la base de datos Chinook}

\begin{center}\small
\begin{tabular}{ll}
\toprule
\textbf{Tabla} & \textbf{Atributos principales} \\
\midrule
\texttt{Artist}       & \texttt{ArtistId (PK), Name} \\
\texttt{Album}        & \texttt{AlbumId (PK), Title, ArtistId (FK)} \\
\texttt{Track}        & \texttt{TrackId (PK), Name, AlbumId (FK), GenreId (FK), MediaTypeId (FK), Milliseconds, UnitPrice} \\
\texttt{Genre}        & \texttt{GenreId (PK), Name} \\
\texttt{Invoice}      & \texttt{InvoiceId (PK), CustomerId (FK), InvoiceDate, BillingCountry, Total} \\
\texttt{InvoiceLine}  & \texttt{InvoiceLineId (PK), InvoiceId (FK), TrackId (FK), UnitPrice, Quantity} \\
\texttt{Customer}     & \texttt{CustomerId (PK), FirstName, LastName, Country, SupportRepId (FK a Employee)} \\
\texttt{Employee}     & \texttt{EmployeeId (PK), FirstName, LastName, ReportsTo (FK a Employee)} \\
\texttt{Playlist}     & \texttt{PlaylistId (PK), Name} \\
\texttt{PlaylistTrack}& \texttt{PlaylistId (FK), TrackId (FK)} \\
\bottomrule
\end{tabular}
\end{center}

\vspace{4pt}
\noindent\rule{\linewidth}{0.4pt}

% ─────────────────────────────────────────────────────────────
\section*{Ejercicio 1 — SQL \hfill\normalfont(0{,}5 puntos)}

Listar el \texttt{Name} de cada género musical junto con la cantidad de canciones que contiene. Incluir también los géneros que no tienen ninguna canción (mostrando 0). Ordenar de mayor a menor cantidad.

\vspace{8pt}

% ─────────────────────────────────────────────────────────────
\section*{Ejercicio 2 — SQL \hfill\normalfont(0{,}75 puntos)}

Listar el \texttt{FirstName} y \texttt{LastName} de los clientes que hayan gastado \textbf{más de \$40 en total} (suma de \texttt{Total} en \texttt{Invoice}). Mostrar también el monto total gastado, ordenado de mayor a menor.

\vspace{8pt}

% ─────────────────────────────────────────────────────────────
\section*{Ejercicio 3 — SQL \hfill\normalfont(1 punto)}

Para cada canción, mostrar su \texttt{Name}, el \texttt{Title} del álbum al que pertenece, su duración en milisegundos (\texttt{Milliseconds}) y la \textbf{duración promedio} de todas las canciones del mismo álbum. Usar una función de ventana con \texttt{AVG(...) OVER (PARTITION BY ...)}. Ordenar por \texttt{Album.Title} y luego por \texttt{Track.Name}.

\vspace{8pt}

% ─────────────────────────────────────────────────────────────
\section*{Ejercicio 4 — Álgebra Relacional \hfill\normalfont(0{,}75 puntos)}

Se tienen las siguientes relaciones del Mundial de Fútbol 2022:

\begin{center}
\begin{tabular}{ll}
\toprule
\textbf{Relación} & \textbf{Atributos} \\
\midrule
\texttt{Players} & \texttt{id, name, birth\_year, playing\_position, national\_team} \\
\texttt{Matches}  & \texttt{id, home, away, datetime, stage} \\
\texttt{Stages}   & \texttt{id, name} \\
\bottomrule
\end{tabular}
\end{center}

\noindent\small\textit{Nota: \texttt{home} y \texttt{away} son códigos de selección (ej.\ \texttt{'ARG'}). \texttt{stage} en \texttt{Matches} referencia a \texttt{Stages.id}.}

\normalsize\vspace{6pt}

Obtener el \texttt{name} y \texttt{national\_team} de todos los jugadores que se desempeñan como delanteros (\texttt{playing\_position = 'FW'}) \textbf{y} que pertenecen a selecciones que jugaron al menos un partido \textbf{como local}. Use las operaciones $\sigma$, $\pi$, $\bowtie$ y eventualmente $\rho$.

\end{document}
